% !TEX root = ../libretto.tex
% Sezione: Natale
\momento{Natale}
\begin{songs}{titoli}

% Fonte: LibrettoAccordi.pdf (Canti Coro), N6.
% Nell'elenco: Comunione, Natale.
\beginsong{Adeste, fideles}[by={Canto popolare}]
\beginverse
Adeste, fideles, laeti triumphantes,
venite, venite in Bethlem:
natum videte Regem angelorum.
\endverse
\beginchorus
Venite, adoremus, venite adoremus,
venite adoremus Dominum.
\endchorus
\beginverse
En grege relicto, humiles ad cunas
vocati pastores adproperant:
et nos ovanti gradu festinemus.
\endverse
\beginverse
Aeterni parentis splendorem aeternum
velatum sub carne videbimus:
Deum infantem pannis involutum.
\endverse
\endsong

% Fonte: LibrettoAccordi.pdf (Canti Coro), E1.
% Nell'elenco: Alleluia, Natale.
\beginsong{Alleluia, è nato Gesù}[by={F. Baggio}]
\beginchorus
Alleluja! Alleluja! Alleluja!
Alleluja! Alleluja! Oggi è nato Gesù!
\endchorus
\beginverse
Al mondo annunciamo la gioia più grande: è nato Cristo Signore.
È l'Emanuele, il Dio con noi, è luce per l'umanità.
\endverse
\beginverse
Un giorno santo è spuntato per noi, venite, adoriamo il Signore.
Da tutta la terra un canto di gioia s'innalza al Dio fatto uomo.
\endverse
\endsong

% Fonte: LibrettoAccordi.pdf (Canti Coro), N3.
% Nell'elenco: Comunione, Natale.
\beginsong{Astro del ciel}[by={Canto popolare}]
\beginverse
Astro del Ciel, pargol divin, mite agnello, Redentor,
tu che i Vati da lungi sognar, tu che angeliche voci nunziar,
\endverse
\beginchorus
Luce dona alle menti, pace infondi nei cuor. \rep{2}
\endchorus
\beginverse
Astro del Ciel, pargol divin, mite agnello, Redentor,
tu di stirpe regale decor, tu virgineo, mistico fior,
\endverse
\beginverse
Astro del Ciel, pargol divin, mite agnello, Redentor,
tu disceso a scontare l'error, tu sol nato a parlare d'amor,
\endverse
\endsong

% Fonte: LibrettoAccordi.pdf (Canti Coro), K11.
% Nell'elenco: Finale, Natale.
\beginsong{Cantate al Signore}[by={D. Ricci}]
\beginverse
Cantate al Signore un canto nuovo da tutta la terra!
Voi cantate al Signore, cantate a lui. Benedite il nome suo.
Annunziate ogni giorno di lui la salvezza e la gloria.
A tutti i popoli dite i prodigi di lui. Benedite il nome suo.
\endverse
\beginchorus
È nato per noi, il Signore per noi oggi è nato.
Lui è nato per noi, il salvatore del mondo è nato.
\endchorus
\beginverse
Esultino i campi, esulti la terra, esultino i mari,
esultino i cieli cantando di lui! Benedite il nome suo.
E sia gioia, perché lui stende già sulla terra la mano sua:
a tutti i popoli la sua giustizia darà. Benedite il nome suo.
\endverse
\endsong

% Fonte: LibrettoAccordi.pdf (Canti Coro), J37.
% Nell'elenco: Comunione, Natale.
\beginsong{Emmanuel}[by={G. Cento}]
\beginchorus
Emmanuel, Emmanuel,
Emmanuel è Dio con noi, è Dio con noi. \rep{2}
\endchorus
\beginverse
Il Dio lontano si è fatto vicino, ha preso un volto come noi.
Lui pensa da uomo, lavora da uomo e ama da uomo: è uno di noi.
\endverse
\beginverse
È nato un bambino, è nato il Signore, è nato il nostro Redentore.
La terra è in festa, un uomo crede, la Chiesa canta: Alleluia!
\endverse
\beginverse
Abbiamo toccato il Verbo di vita, abbiamo parlato insieme a lui.
Ha dato la vita per l'uomo che crede, per tutta la nostra umanità.
\endverse
\endsong

% Fonte: LibrettoAccordi.pdf (Canti Coro), N2.
% Nell'elenco: Finale, Natale.
\beginsong{Gli angeli delle campagne}[by={Canto popolare}]
\beginverse
Gli angeli delle campagne, cantano l'inno: Gloria in ciel!
E l'eco delle montagne ripete il canto dei fedel.
\endverse
\beginchorus
Gloria in excelsis Deo! \rep{2}
\endchorus
\beginverse
Oh pastori che cantate, dite il perché di tanto onor.
Qual Signore, qual Profeta merita questo gran splendor?
\endverse
\beginverse
Oggi è nato in una stalla, nella notturna oscurità.
Egli è il Verbo, s'è incarnato, e venne in questa povertà.
\endverse
\endsong

% Fonte: LibrettoAccordi.pdf (Canti Coro), N5.
% Nell'elenco: Comunione, Natale.
\beginsong{Lieto annuncio}[by={Canto popolare}]
\beginverse
Lieto annuncio risuonò, luce nuova sfolgorò: Cristo è nato all'umanità.
\endverse
\beginchorus
Hai vinto ormai le tenebre: sei nato Re del mondo. \rep{2}
\endchorus
\beginverse
Gioia piena brilla in noi; vive qui l'eternità; nasce nuova umanità.
\endverse
\endsong

% Fonte: LibrettoAccordi.pdf (Canti Coro), G8.
% Nell'elenco: Offertorio, Natale.
\beginsong{Nasce l'amor}[by={Gen Rosso}]
\beginverse
Fredda è la notte, buio d'intorno, silente attesa, sacro timor.
Brilla una stella, veglia la terra, nasce la vita, palpita un cuor.
\endverse
\beginchorus
Nasce Gesù, nasce l'Amor. Gli angeli in cielo cantano in cor. \rep{2}
\endchorus
\beginverse
E se l'amore regna fra noi, nei nostri cuori vive Gesù.
Splende la luce, in tutti è pace: il paradiso è in mezzo a noi.
\endverse
\beginverse*
Nato è Gesù, nato è l'Amor. Gli angeli in cielo cantano in cor. \rep{2}
\endverse
\endsong

% Fonte: LibrettoAccordi.pdf (Canti Coro), K9.
% Nell'elenco: Finale, Natale.
\beginsong{Oggi è nato}[by={D. Ricci}]
\beginverse
Oggi per noi è nato il Salvatore. Tutta la terra canta a lui, benedice lui.
E ogni giorno si leva un canto, canto nuovo di gioia e vita.
E ogni giorno si narra la gloria ed osannano i popoli che oggi è nato.
\endverse
\beginchorus
Tutto il mondo canta alleluja. Anche il cielo canta alleluja.
Canta la terra, canta il mare
e nei campi le messi rispondono agli alberi. \rep{2}
\endchorus
\beginverse*
Tutto il mondo canta alleluja. Anche il cielo canta alleluja.
Canta la terra, canta il mare,
alleluja, alleluja, alleluja, alleluja.
\endverse
\endsong

% Fonte: LibrettoAccordi.pdf (Canti Coro), N1. Spazi tra le parole latine ricostruiti.
% Nell'elenco: Comunione, Natale (canto popolare).
\beginsong{Resonet in laudibus}[by={Canto popolare}]
\beginverse
Resonet in laudibus cum iucundis plausibus Sion cum fidelibus.
\endverse
\beginchorus
Apparuit, apparuit quem genuit Maria. \rep{2}
\endchorus
\beginverse
Sion lauda Dominum, Salvatorem omnium: Virgo parit Filium.
\endverse
\beginverse
Pueri concurrite, nato Regi psallite, voce pia dicite.
\endverse
\beginverse
Ergo nostra concio cum sit plena gaudio benedicat Domino.
\endverse
\beginverse
Et nos unanimiter proclamemus dulciter ipse Summus Arbiter.
\endverse
\endsong

% Versione popolare (3 strofe). Fonte: spartito Singing Bell,
% singing-bell.com/wp-content/uploads/2015/11/Tu-scendi-dalle-stelle_Singing-Bell.pdf
% Testo originale completo (7 strofe): it.wikipedia.org/wiki/Tu_scendi_dalle_stelle
% Nell'elenco: Comunione, Natale.
\beginsong{Tu scendi dalle stelle}[by={S. Alfonso Maria de' Liguori}]
\beginverse
Tu scendi dalle stelle, o Re del cielo,
e vieni in una grotta al freddo e al gelo. \rep{2}
O Bambino mio divino,
io ti vedo qui a tremar;
o Dio beato!
Ah, quanto ti costò l'avermi amato! \rep{2}
\endverse
\beginverse
A te, che sei del mondo il Creatore,
mancano panni e fuoco, o mio Signore. \rep{2}
Caro eletto pargoletto,
quanto questa povertà
più m'innamora,
giacché ti fece amor povero ancora. \rep{2}
\endverse
\beginverse
Tu lasci il bel gioir del divin seno,
per giungere a penar su questo fieno. \rep{2}
Dolce amore del mio cuore,
dove amor ti trasportò?
O Gesù mio,
perché tanto patir? Per amor mio! \rep{2}
\endverse
\endsong

\end{songs}
